\chapter{Grupo motor}
\label{cap:mm-motor}

\opening{Un motor de 180 caballos limitado a 124: el margen está en las placas de presión de admisión, no en el acelerador.}

\section{Descripción}
\fuente{MM, capítulo 6}

Este capítulo cubre el retiro y la instalación del motor y de sus sistemas de soporte (inducción, enfriamiento y escape). Para la inspección, reparación y mantenimiento específicos del motor y sus componentes, consulte las publicaciones de sus fabricantes.

Los R22 usan uno de cuatro motores Lycoming:

\begin{itemize}
\item \textbf{O-320-A2B u O-320-A2C:} R22 Standard; potencia nominal 150~HP.
\item \textbf{O-320-B2C:} R22 HP, Alpha y Beta; 160~HP.
\item \textbf{O-360-J2A:} Beta II y Mariner II; capaz de 180~HP, reducido por Lycoming a 145~HP.
\end{itemize}

Todos los motores están limitados a \textbf{124~HP de potencia máxima continua} mediante la limitación de presión de admisión y RPM. Beta, Mariner y Beta II/Mariner II tienen una potencia de despegue de 5 minutos de 131~HP.

El arranque del O-320-A2B se asiste con un acoplamiento de impulso en el magneto izquierdo del motor (derecho del helicóptero). En el O-320-A2C, el O-320-B2C y el O-360-J2A se asiste con un sistema vibrador de arranque y platinos de retardo en el magneto izquierdo del motor. Los platinos del tacómetro, en el magneto derecho del motor, entregan la RPM al controlador del gobernador y al tacómetro. Con 104\,\% en el tacómetro, el motor gira a 2652~RPM.

La inducción pasa por un filtro de aire en una caja de aire montada en el carburador. La calefacción la provee una toma montada en el tubo de escape y un ducto, y se controla con una válvula de guillotina accionada por cable en la caja de aire. \textbf{El aire de entrada al carburador está filtrado incluso con la calefacción conectada.} Ante un filtro obstruido, una válvula de derivación con resorte en la tapa de la caja de aire, o la flexión de la caja, deja pasar el aire. En el Beta II/Mariner II la calefacción se correlaciona con el colectivo mediante un embrague de fricción (asistente de calefacción), con un pestillo para trabarla cerrada.

El enfriamiento lo da un ventilador montado en el motor, encerrado en una carcasa de fibra de vidrio que conduce el aire a paneles montados en el motor, que a su vez lo dirigen a los cilindros, al enfriador de aceite externo, al alternador y a la caja principal. La potencia se transmite por un sistema vertical de poleas y correas acopladas por un actuador eléctrico de tensión. La carga de batería y la energía eléctrica las provee un alternador de 60~A accionado por correa (opcional de 70~A).

\section{Retiro e instalación del motor}

\begin{nota}
Si el motor se retira para reemplazar un bastidor inferior, cambiar soportes o reparar el cortafuego, el ventilador puede quedar instalado; si el motor se va a desarmar, retire el ventilador.
\end{nota}

\textbf{Retiro, puntos clave:} desconecte la línea estática y retire el capó trasero, el carenado del cono de cola y el bastidor B429-5 (si hay tanque auxiliar, desabastezca y retírelo); desconecte el ducto de enfriamiento de la caja principal y retire el panel de acceso del cortafuego horizontal; retire el embrague (\S\,7.210); cierre el combustible y desconecte y tape la manguera del carburador; desconecte el acelerador, la mezcla, la calefacción y la sonda de CAT (conector rápido a unas 4~in de la sonda); retire carburador y caja de aire (pueden salir juntos); desconecte y tape la manguera de presión de aceite y la línea de MAP del cilindro delantero izquierdo; identifique y desconecte masas, alternador, arranque, sensores y el motor del actuador; ducto de calefacción y cebador si existen. Ponga acolchado protector, enganche el útil MT016-1 a un izaje de \textbf{500~lb mínimo} y al gancho de izaje del motor. \textbf{Anote los suplementos de cada soporte del motor}, retire los cuatro pernos de montaje (hay que retirar la tapa de bujías derecha), apoye el patín de cola o retire el cono, retire el montante trasero derecho y saque el motor entre dos personas, variando su altura para librar el tubo transversal trasero.

\textbf{Instalación, puntos clave:}
\begin{itemize}
\item Use los suplementos anotados. Si no hay registro, ponga \textbf{0,063~in} entre ambos soportes inferiores y el motor, y \textbf{0,190~in} entre ambos superiores. Torque y franja; instale el montante trasero derecho.
\item Carburador con junta nueva y torque especial; \emph{palnuts} y franja.
\item Acelerador al \textbf{agujero central} del balancín del carburador (reglaje en \S\,10.150). Mezcla y calefacción con \textbf{1/16~in de retroceso} en cada extremo del recorrido.
\item Masa del motor y pernos del cárter con su torque especial; líneas de aceite y de MAP; conexiones eléctricas; ductos y abrazaderas; combustible.
\item Antes de hacer funcionar el motor, verifique: alineación de poleas (\S\,7.230); suplementado del eje de cola (\S\,7.330); holgura carenado de enfriamiento--cono de cola \textbf{0,250~in mínimo}; holgura placa flexible intermedia--balancín del rotor de cola \textbf{0,250~in mínimo con correas de más de 10~h, 0,125~in con correas nuevas}; ángulo del eje del embrague (\S\,7.240) y centrado de la polea superior (\S\,7.250). Llene el aceite.
\end{itemize}

\begin{advertencia}
Un suplementado incorrecto del motor puede dañar o hacer fallar componentes de la transmisión.
\end{advertencia}

\section{Suplementado del motor}

Para una correcta alineación de las poleas, el motor debe suplementarse de modo que la distancia entre la parte superior de la corona dentada y la parte inferior del tubo transversal largo del bastidor superior sea de \textbf{3,55 a 3,75~in con el embrague acoplado}. (Puede fabricarse un calibre pasa/no pasa de chapa rígida de 3,55 × 3,75~in con un agujero central de 2~in.)

\begin{pasos}
\item Acople el embrague y mida. Si excede 3,75~in, agregue suplementos iguales entre los bastidores inferiores y los soportes de goma inferior izquierdo e inferior derecho. Espesor máximo calculado:
\[
\text{suplemento} = \frac{\text{medida} - 3{,}55}{2}
\]
Elija suplementos hasta ese valor sin superarlo. Hay suplementos de 0,032, 0,063 y 0,190~in. Quitar suplementos aumenta la distancia.
\item Desacople el embrague y retire los faldones.
\item Afloje los pernos de los soportes izquierdos (dejando apretados los derechos), instale los suplementos entre el bastidor inferior izquierdo y su soporte y aplique torque. Luego repita en el lado derecho (retirando la tapa de bujías derecha) dejando apretados los izquierdos.
\item Acople el embrague, verifique la medida, ajuste y marque las tuercas con franja.
\item Verifique la alineación de poleas (\S\,7.230), reinstale tapas y faldones y ajuste la correlación del acelerador (\S\,10.150).
\end{pasos}

Suplementos: arandelas A141-12 (0,190~in), A141-24 (0,032~in), A141-25 (0,063~in) y B155 (0,032~in). Debe haber al menos una A141-12 o -25 junto al soporte de goma; las B155 van entre dos arandelas A141. La B155-2 encaja a presión (sin chaveta ni alambre); la B155-1 requiere chaveta o alambre de frenado. Use arandelas para dejar 2 a 4 hilos expuestos.

\begin{nota}
Los suplementos agregados o quitados en los dos soportes inferiores cambian la distancia corona--bastidor aproximadamente el doble de su espesor: agregar 0,032~in en cada soporte inferior eleva la parte trasera del motor unos 0,064~in. Los pernos de soporte AN6 son obsoletos y deben reemplazarse por NAS1306 o NAS6606 al suplementar.
\end{nota}

\section{Sistema de enfriamiento}

\subsection{Retiro del ventilador}

\begin{precaucion}
No acople el actuador del embrague sin la carcasa instalada.
\end{precaucion}

Si se reinstalará el mismo ventilador, oriente el soporte de la corona en la posición de sincronización del magneto y marque el ventilador en la línea de partición de la carcasa, para minimizar la necesidad de nuevos pesos de balance.

\begin{pasos}
\item Desacople por completo el actuador, abra el disyuntor \textsf{CLUTCH START} (5 o 7,5~A) y apague la batería.
\item Desconecte la manguera de calefacción de la mitad inferior de la carcasa, si existe. Retire los tornillos perimetrales carcasa--paneles y los dos pernos que unen la carcasa al rodamiento A181-4.
\item Corte el alambre y retire el pasador elástico de la tuerca del ventilador; retire tuerca y arandelas y vuelva a enroscar la tuerca 3 vueltas completas para que el ventilador no caiga.
\item Instale los espaciadores MT092-3 en los tres pernos largos B660-1 (o reemplace tres tuercas alternas por espaciadores), lubrique el perno central del extractor MT092-1 con antiagarrotante A257-9 y móntelo. Inmovilice la corona con el útil MT091-1.
\item Apriete el perno central del extractor hasta que el ventilador se suelte del eje cónico (puede salir con un fuerte estallido).
\item Retire el extractor y aplique protector B270-21 a las roscas expuestas de los B660-1. Si se retiraron tuercas, aplique imprimante húmedo, reinstale arandelas y tuercas y, con el imprimante húmedo, dé el torque especial y la franja.
\item Retire la tuerca, el ventilador y la carcasa.
\end{pasos}

\subsection{Instalación del ventilador}

\begin{precaucion}
El Beta II/Mariner II debe usar ventiladores B174-1 rev.~F o posterior (12 pernos NAS6603-6 de unión cono--ventilador) y ejes A007-5 o -6 (tuerca de 1-11/16~in) o posteriores.
\end{precaucion}

\begin{pasos}
\item Evalúe el estado de las correas, la correa del alternador y el rodamiento A181 y reemplácelos si es necesario. Limpie y seque con solvente aprobado el eje cónico y el cubo del ventilador; inspecciónelos.
\item Coloque la mitad superior de la carcasa sobre el ventilador, instale el ventilador en el eje y fije la carcasa a los paneles.
\item Con antiagarrotante A257-9 \emph{solo} en la parte roscada del eje A007 y en las caras de apriete de la arandela MS20002-18 (contra el cubo), las arandelas NAS1149F1832P y la tuerca almenada, instale la ferretería. (Si es el ventilador original, ubique el soporte de la corona en la posición de sincronización y alinee las marcas.)
\item Con la corona inmovilizada (MT091-1), dé el torque especial a la tuerca hasta que su ranura coincida con el agujero del eje; instale el pasador elástico y fréne con alambre inoxidable de 0,041~in (puede hacerse después del balanceo). Si no coincide en todo el rango de torque, agregue o quite arandelas bajo la tuerca.
\item Verifique el torque de los seis pernos del cubo del ventilador (franja después del balanceo).
\item Instale la mitad inferior de la carcasa e introduzca los dos pernos AN3 por el soporte inferior y el bloque del rodamiento del actuador hasta las tuercas del soporte superior (alinee con un destornillador en la muesca); torque estándar y franja. Si se reemplazó un ventilador rev.~I o anterior por uno rev.~J o posterior sin cambiar la carcasa inferior, consulte a RHC para reubicar la tapa D229-4.
\item Instale los tornillos perimetrales. Marque el ventilador con sellador de torque frente a cada extremo del pasador de la tuerca para detectar movimiento entre ventilador y eje.
\item Tire del ventilador hacia atrás para absorber el juego del cigüeñal y verifique una holgura labio B229--entrada del ventilador de \textbf{0,060 $\pm$ 0,030~in} (pueden ovalarse los agujeros del labio).
\item Balancee dinámicamente.
\end{pasos}

\begin{precaucion}
El balance del ventilador debe verificarse al instalarlo: el desbalance puede causar daños. La tuerca del ventilador debe reapretarse después del primer funcionamiento del motor.
\end{precaucion}

\subsection{Balance dinámico del ventilador}

Instrucciones para Chadwick-Helmuth Vibrex 2000 (para otros equipos, siga sus instrucciones):

\begin{pasos}
\item Lave el ventilador y sus paletas con jabón suave (pH 7 a 9) y agua tibia; enjuague con agua potable.
\item Instale el acelerómetro en el perno izquierdo del rodamiento inferior del actuador, con el cable hacia abajo; la fotocelda lo más cerca posible de las 12 en punto en la carcasa superior, y la cinta reflectante C701-1 o 907-4270-1 en el borde de una paleta.
\item Haga funcionar en tierra a 104\,\% con el rotor acoplado (\textbf{extreme el cuidado cerca del rotor de cola}); con el motor estabilizado en unas 2652~RPM, mida y anote la vibración y el ángulo de reloj. Detenga el motor.
\item Tras el primer funcionamiento, reapriete la tuerca AN320-18 con el torque especial e instale y frene el pasador; verifique el torque de los seis pernos del cubo.
\item \textbf{Si la vibración supera 0,2~IPS, cambie los pesos.} Con la cinta en la posición anotada, aumente el peso bajo la tuerca más cercana a la 1 en punto (o reparta entre las dos más cercanas), o disminúyalo bajo la más cercana a las 7. Efecto aproximado de cada arandela: NAS1149F0363P, 0,03--0,05~IPS; A141-17, 0,08--0,10~IPS; AN970-3, 0,18--0,20~IPS (pueden combinarse). Mantenga 2 a 4 hilos expuestos; no retire los pernos NAS6603-6 (cuide las arandelas NAS1149F0316P); dé el torque especial. El acceso es por el agujero bajo la tapa D229-4 (kit KI-111) o por la compuerta del enfriador de aceite. Radio de los pesos: 5,10~in en ventiladores rev.~J y posteriores; 4,10~in en rev.~I y anteriores.
\item Repita hasta lograr \textbf{menos de 0,2~IPS}; retire el equipo y aplique las franjas.
\end{pasos}

\begin{precaucion}
Asegúrese de que el encendido esté apagado y sin llave siempre que mueva el ventilador a mano.
\end{precaucion}

\subsection{Reemplazo de la carcasa de enfriamiento}

Nivele lateralmente la célula; arme los soportes A185-1 (abajo), A185-2 (arriba izquierda) y A185-5 (arriba derecha) en el alojamiento del rodamiento y nivele el rodamiento inferior del actuador. Monte la mitad superior y el ventilador (en Beta y Mariner con enfriador grande, su lado derecho apoya en el reborde del enfriador), centre el labio del ventilador con el borde de la toma, agregue la mitad inferior y verifique una holgura concéntrica. Marque y taladre (broca \#18) los agujeros de montaje a través de los paneles, ubique los labios B229-2 a \textbf{0,030--0,070~in} del ventilador y taládrelos; taladre los soportes y placas de refuerzo (\#40, luego \#30) y remache con MS20470AD4. Instale una tuerca remachada NAS679A08 en el agujero vertical interior derecho, recorte los bordes para dejar 0,01--0,02~in de distancia a las pestañas de los paneles y sin contacto con la cabeza del perno del actuador; imprima con epóxico y pinte de gris.

\section{Sistema de lubricación}

Los grados y la cantidad de aceite están en el capítulo~\ref{cap:manejo}; para su ajuste y mantenimiento, vea el manual del operador de Lycoming.

\textbf{Enfriador de aceite:} para retirarlo, desconecte las líneas, los cuatro tornillos del soporte a los paneles y los cuatro pernos de montaje. Al reemplazarlo, \textbf{no intente retirar ni transferir los codos de 90° anodizados en azul} (los de acero sí pueden retirarse). Selle las roscas cónicas de los codos con sellador B270-6, sujete con llave los acoples roscados del enfriador al apretar los codos, apriete los pernos AN3 y las tuercas de las líneas con su torque y marque las franjas.

\section{Sistema de inducción}

\subsection{Carburador y caja de aire}

Para retirar el carburador: desconecte las mangueras de aire y calefacción, el acelerador, la mezcla, el cable de calefacción y la sonda de CAT, y retire las \emph{palnuts} y tuercas del múltiple (sale con la caja de aire).

\textbf{Instalación de la caja de aire:}
\begin{pasos}
\item Con las tomas de aire a la derecha, instale la caja y su junta de goma con pernos y arandelas (o tornillos de cabeza y espaciadores en los S/N 2571 a 2664 con caja C057); apriete y frene con alambre de 0,032~in.
\item Conecte el resorte de seguridad de la mezcla al brazo del carburador y verifique que lleve el brazo a rica a fondo.
\item Conecte la funda del cable de mezcla a la caja (empuje y tire: no debe resbalar) y el cable interior al brazo: el control debe llevar el brazo de tope a tope, con \textbf{0,03 a 0,10~in de retroceso} bajo el pomo en rica a fondo. Apriete el acople A462, \emph{palnut} y franja.
\item Conecte el cable de calefacción a la válvula deslizante: debe abrirla y cerrarla por completo. Doble a 90° el extremo libre del cable interior al salir del acople y córtelo dejando 0,25 a 0,50~in.
\item Instale el filtro y cierre la tapa apretando los cuatro pernos largos solo lo necesario para eliminar el juego; franja. En los Mariner, selle el perímetro de la tapa con cinta de aluminio. Verifique que la compuerta alternativa con resorte cierre por completo.
\item Conecte las mangueras de entrada.
\end{pasos}

\textbf{Instalación del carburador:} con el balancín de la mariposa a la izquierda, instálelo con junta nueva sobre los espárragos del cárter con torque especial, \emph{palnuts} con torque estándar y franja; el tubo del acelerador al agujero central del balancín; conecte el cable de mezcla sin dar torque al acople (se aprieta al instalar la caja de aire); instale la caja; verifique ralentí y mezcla de ralentí (\S\,2.210) y la correlación (\S\,10.150).

\subsection{Limpieza del filtro de aire}

\begin{nota}
Reemplace los filtros de aire B771-1 a los 5 años o 600 horas de vuelo, lo que ocurra primero.
\end{nota}

\begin{pasos}
\item Retire el filtro e inspecciónelo. Si el medio plisado solo tiene polvo, límpielo con aire comprimido o agua; si está muy sucio, con jabón suave y agua:
\begin{itemize}
\item \textbf{Aire comprimido:} boquilla al menos a 1~in del medio, menos de 40~psi, en sentido contrario al flujo normal.
\item \textbf{Agua:} chorro a menos de 40~psi en sentido contrario al flujo; seque a fondo con aire tibio a menos de 160~°F.
\item \textbf{Jabón y agua:} chorro a menos de 40~psi en sentido contrario; remoje en solución jabonosa suave más de 15 minutos y menos de 4 horas, agitando suavemente; enjuague con chorro a menos de 40~psi en ambos sentidos; seque con aire a menos de 160~°F.
\end{itemize}
\item Con una luz brillante, verifique que el medio no tenga agujeros ni roturas.
\item Verifique que las superficies de sellado sean un círculo continuo, liso y plano.
\item Instale el filtro. \textbf{No instale un filtro húmedo ni aplique aceite al medio.}
\end{pasos}

\subsection{Toma de aire caliente}

\textbf{O-320:} la toma va sobre una abrazadera de cordón en el tubo vertical del escape del cilindro N.º~4, con una abrazadera A694 y un ángulo, y se cubre con una chapa curva de seis tornillos. Al instalarla, mueva la abrazadera y gire la de cordón para minimizar la precarga y asegurar holgura con la estructura; pernos con torque estándar. \textbf{O-360:} la toma va con dos abrazaderas B277 al múltiple de escape, con torque especial de \textbf{30~in-lb}, posicionada para minimizar la precarga.

\section{Sistema de escape}

Para retirarlo: retire la manguera de calefacción del carburador y las de calefacción de cabina, proteja el tubo transversal trasero, desconecte las bridas en las culatas y retire el escape y las juntas.

Para instalarlo: junta en cada lumbrera; proteja el tubo transversal; conecte las bridas con arandelas y tuercas; afloje la abrazadera de cordón de cada tubo vertical; dé el torque especial a las tuercas de las bridas, instale \emph{palnuts} con torque estándar; reinstale las mangueras; apriete los pernos AN3 de las abrazaderas con torque estándar y reposicione para minimizar la precarga.

\begin{nota}
No use cinta inoxidable bajo las abrazaderas de cordón del escape que incorporan barrera integral.
\end{nota}

\section{Diagnóstico de fallas}

\subsection{Lista de chequeo por baja potencia}

Cuando se sospeche baja potencia del motor, verifique:

\begin{pasos}
\item Que no se exceda el peso bruto (pese la aeronave si no se pesó en el último overhaul).
\item Presión y temperatura de aceite y CHT dentro de límites.
\item El indicador de MAP marca la presión ambiente con el motor detenido.
\item Grado de combustible correcto.
\item Palas del rotor principal limpias y lisas (sin pintura áspera).
\item Escape sin obstrucciones.
\item Aletas compensadoras (\emph{trim tabs}) de las palas sin dobleces excesivos.
\item La válvula de calefacción cierra por completo en OFF.
\item Inducción sin obstrucciones: filtro y ambas mangueras (ambiente y calefacción) limpios, sanos y que no colapsen en vuelo; junta carburador--motor del tamaño correcto.
\item Indicación correcta del tacómetro del motor (verifique con un tacómetro digital, como el de los equipos de balanceo modernos).
\item Mallas de aceite o filtro sin residuos metálicos indicativos de daño interno.
\item Compresión de cilindros correcta (fuga y directa) según Lycoming.
\item Mezcla ni excesivamente pobre ni rica: compare el flujo de combustible, revise la limpieza del filtro de combustible, las bujías y la sujeción de los tornillos cuerpo del acelerador--cuba (Lycoming SB 366).
\item Tipo y estado correctos de las bujías.
\item Sincronización magneto--motor correcta. Verifique que el magneto izquierdo del motor (derecho del helicóptero) no opere en modo retardo con la llave en BOTH: desconecte el terminal de retardo tras arrancar y compare el MAP en estacionario; si con el terminal desconectado el MAP es menor, el magneto operaba retardado.
\item Flujo de combustible correcto y malla de entrada del carburador sin obstrucciones (prueba de flujo del capítulo 12 del MM).
\item Con cebador, que no haya fuga de combustible por la bomba (trabada) ni de aire por los acoples de la culata.
\item El brazo del acelerador del carburador toca el tope de apertura total (colectivo levemente arriba).
\item El brazo de la mezcla toca su tope en rica a fondo y, desconectado del cable y del resorte, no muestra retroceso en los extremos (agarrotamiento interno).
\item Sin fugas en la inducción ni en el sistema de MAP. Para el MAP: desconecte la línea en la culata, aplique succión con una jeringa hasta 10~inHg y observe un minuto: \textbf{no debe subir más de 1~inHg}. Verifique que los bujes del eje del acelerador no estén gastados (holguras axial y radial). También puede haber fugas por conexiones flojas de los tubos de admisión en el cárter o tubos agrietados. Solo en el O-320-B2C: los tubos de admisión de brida delgada P/N 61363, 61364 o 61365 usados en lugar de los correctos LW-16728, LS-16729 y LW-16730 pueden causar fugas con la brida de contrataladro profundo P/N 74360.
\item Holgura de las guías de válvulas de escape correcta (Lycoming SB 388B).
\item Holgura en seco de los taqués correcta (manual de overhaul de Lycoming).
\item Con los taqués inflados (sin juego), cada válvula de admisión debe tener igual carrera de apertura medida con un comparador, y lo mismo las de escape; una carrera desigual indica levas dañadas.
\item Sincronización interna correcta de los magnetos.
\item Nivel de combustible correcto en la cuba del carburador.
\item La lumbrera de admisión de la culata tiene un borde redondeado donde el flujo gira 90° hacia la válvula; si no, consulte a Lycoming.
\item Sincronización cigüeñal--árbol de levas correcta: cuando el pistón N.º~1 pasa el PMS en compresión, la válvula de admisión N.º~2 debe estar abriendo mientras la de escape N.º~2 cierra (traslapo).
\item Si un O-320/O-360/O-540 tiene baja potencia justo después de un overhaul, verifique que se instalaron pistones de relación 8,5:1 P/N 75089 (y no los de baja compresión 7,0:1 P/N 75413).
\end{pasos}

\subsection{Temperatura de culata alta}

Ante CHT inusualmente altas, verifique: grado de combustible correcto; holgura labios metálicos de la carcasa--ventilador de \textbf{0,020 a 0,100~in}; mezcla en rica a fondo en el carburador con el pomo completamente adentro; integridad de los deflectores y paneles de enfriamiento; mangueras de enfriamiento bien ruteadas, aseguradas y en buen estado; bujías correctas (Lycoming SI 1042); magnetos bien sincronizados; sin fugas excesivas en las juntas tubo de escape--culata. Las CHT pueden ser más altas durante el asentamiento, hasta el cambio de aceite mineral a dispersante sin cenizas. Si nada mejora, el indicador de CHT puede estar marcando erróneamente alto (\S\,33-131-2 parte E).

\section{Control del acelerador y asistente de calefacción}

El control del acelerador consta de la empuñadura giratoria integrada al colectivo, el resorte de sobrerrecorrido, el balancín, el tubo de empuje--tracción, la palanca del acelerador del carburador y, desde el S/N 2571, el asistente de calefacción instalado en fábrica.

\textbf{Resorte de sobrerrecorrido A327-1:} se accede retirando el panel del vientre y el panel vertical entre los respaldos. Al instalarlo, conecte primero el terminal inferior al balancín y luego el superior al brazo del acelerador del colectivo; con el colectivo arriba y la empuñadura completamente cerrada, verifique la dimensión de la figura 10-8 del MM y ajuste la posición del terminal superior. Torque, \emph{palnuts} y franjas.

\textbf{Tubo del acelerador:} se pasa hacia adelante por el sello del cortafuego, se une al agujero central de la palanca del carburador y luego al balancín bajo el compartimiento de equipaje izquierdo; ajuste la correlación (\S\,10.150) antes de apretar contratuercas y \emph{palnuts}.

\textbf{Cable de calefacción con asistente:} lubrique el cable interior con grasa A257-1. Con el pomo completamente abajo, verifique \textbf{0,10 a 0,25~in} entre el acople del cable interior y el extremo delantero de la funda; el extremo trasero de la funda debe sobresalir \textbf{0,00 a 0,15~in} dentro de la caja de aire; con el pomo abajo y la válvula completamente atrás, corte el cable interior \textbf{0,03 a 0,10~in} más allá del acople. Reponga el sellador B270-5 en el paso del cortafuego.

\textbf{Balancín D334-5 del asistente:} los tornillos A142 del bloque de rodamiento se aprietan a \textbf{45~in-lb}. La \textbf{fricción rotacional debe ser de 10 a 12~in-lb}, medida sujetando el lado largo de la placa D333-11 en una morsa acolchada y girando la cabeza del perno con una llave de torque con indicador; se ajusta con muñones B250 de distinto largo (B250-5: 0,059~in; -6: 0,062; -7: 0,065; -10: 0,068; -12: 0,071), que pueden mezclarse: \textbf{un muñón más delgado aumenta la fricción}. El perno NAS6606 se aprieta a \textbf{280~in-lb} (incluye el torque de autobloqueo de la tuerca), con \emph{palnut} y franja.
